% ============================================================
% Section 112: 泡利自旋磁化率与温度依赖
% ============================================================
\section{固体物理：泡利自旋磁化率与温度依赖}
\label{sec:pauli-susceptibility}

\begin{modelbox}[题目：独立电子气的自旋磁化率]
(1) 证明``独立''电子气在绝对零度时的自旋磁化率为
\begin{equation*}
\chi=\mu_B^2\,g(E_F),
\end{equation*}
其中 $\mu_B$ 是玻尔磁子，$g(E_F)$ 是费米能级处的态密度；

(2) 计算自由电子气（电子--离子相互作用忽略）的 $g(E_F)$；

(3) 自旋磁化率微弱地依赖于温度的形式
\begin{equation*}
\chi(T)=\chi(0)\left(1+aT^2\right),
\end{equation*}
你认为 $a$ 是正的还是负的？（定性论证即可）。
\end{modelbox}

\begin{tipbox}[考点快速分析]
\begin{itemize}
    \item \textbf{泡利顺磁图像}：磁场使自旋向上/向下子能带发生相对位移，电子翻转导致净磁矩
    \item \textbf{态密度}：3D 自由电子气 $g(E_F)=3N/(2E_F)$
    \item \textbf{温度修正}：索末菲展开 $E_F(T)=E_F(0)[1-\pi^2(k_BT/E_F)^2/12+\cdots]$
\end{itemize}
\end{tipbox}

\begin{derivationbox}[标准解题过程]
\textbf{(1) 绝对零度泡利顺磁化率的证明}

无外磁场时，自旋向上和向下的电子填充到相同费米能级 $E_F$，总磁矩为零。

施加磁场 $B$ 后，自旋与磁场平行（磁矩向上）的电子能量降低 $\mu_BB$；反平行（磁矩向下）的电子能量升高 $\mu_BB$。两子能带费米能级相对位移，自旋反平行的电子翻转为平行，直到两者费米能级再次对齐。

翻转的电子数等于费米面附近能量宽度 $\mu_BB$ 范围内的状态数的一半（态密度 $g(E_F)$ 包含两种自旋）：
\begin{equation}
\Delta N\approx\frac{1}{2}\,g(E_F)\cdot(\mu_BB).
\end{equation}

每个翻转电子，磁矩由 $-\mu_B$ 变为 $+\mu_B$，改变量为 $2\mu_B$。总磁矩
\begin{equation}
M=\Delta N\cdot2\mu_B=\frac{1}{2}g(E_F)\mu_BB\cdot2\mu_B=\mu_B^2\,g(E_F)\,B.
\end{equation}

磁化率 $\chi=M/B$，故
\begin{equation}
\boxed{\chi=\mu_B^2\,g(E_F).}
\end{equation}
（SI 单位制中需乘以 $\mu_0$：$\chi=\mu_0\mu_B^2g(E_F)$；高斯制中 $\mu_0=1$。）

\textbf{(2) 自由电子气的态密度 $g(E_F)$}

三维自由电子气（包含自旋简并）的态密度：
\begin{equation}
g(E)=\frac{V}{2\pi^2}\left(\frac{2m}{\hbar^2}\right)^{3/2}E^{1/2}.
\end{equation}

在费米能级处，利用 $E_F=\hbar^2(3\pi^2n)^{2/3}/(2m)$（$n=N/V$），化简得
\begin{equation}
g(E_F)=\frac{3N}{2E_F}.
\end{equation}

因此泡利磁化率也可写为
\begin{equation}
\chi=\frac{3}{2}\frac{N\mu_B^2}{E_F},
\end{equation}
与题 3.9 公式一致。

\textbf{(3) 温度依赖系数 $a$ 的符号}

有限温度下，费米--狄拉克分布使费米面附近电子被热激发。由于态密度 $g(E)\propto E^{1/2}$ 随能量增加而增大，电子从低于 $E_F$ 的态跃迁到高于 $E_F$ 的态时，为保持总电子数守恒，费米能级必须略微\textbf{降低}。

由索末菲展开：
\begin{equation}
E_F(T)=E_F(0)\left[1-\frac{\pi^2}{12}\left(\frac{k_BT}{E_F(0)}\right)^2+\cdots\right].
\end{equation}

泡利磁化率 $\chi\propto g(E_F)\propto E_F^{1/2}$。由于 $E_F(T)<E_F(0)$，故 $g(E_F(T))<g(E_F(0))$，磁化率随温度升高而\textbf{减小}。

展开到 $T^2$：
\begin{equation}
\chi(T)\approx\chi(0)\left[1-\frac{\pi^2}{24}\left(\frac{k_BT}{E_F(0)}\right)^2\right].
\end{equation}
因此系数 $a$ 为\textbf{负值}（$a<0$）。
\end{derivationbox}

\begin{formulabox}[答案汇总]
\begin{center}
\begin{tabular}{ll}
\toprule
\textbf{项目} & \textbf{结果} \\
\midrule
(1) 泡利磁化率 & $\chi=\mu_B^2\,g(E_F)$（高斯制） \\
(2) 态密度 & $g(E_F)=3N/(2E_F)$ \\
(3) 温度系数 & $a<0$（负值），因 $E_F(T)$ 随 $T$ 升高下降 \\
\bottomrule
\end{tabular}
\end{center}
\end{formulabox}

\begin{insightbox}[物理图像]
\begin{itemize}
    \item 泡利顺磁性是金属自由电子的基本磁性，源于费米面附近电子的自旋翻转。其磁化率远小于经典居里顺磁性，且几乎与温度无关。
    \item 自由电子气的态密度随能量以 $E^{1/2}$ 增长，这是导致 $a<0$ 的根本原因。若态密度随能量减小，则 $a$ 可能为正。
    \item 金属的实测磁化率还包含朗道抗磁性（轨道效应）、电子--电子相互作用增强因子等，但泡利顺磁是其主导项之一。
\end{itemize}
\end{insightbox}
